\subsection{完备有界集与拟完备空间}

定义 6. --- 称局部凸空间 E 是拟完备的，如果 E 的每个闭且有界的子集都是完备的（对由 E 的一致结构诱导的一致结构而言）。

完备的局部凸空间是拟完备的，但其逆并不总成立。* 例如，若 E 是一个无穷维 Hilbert 空间，或更一般地，是一个无穷维自反 Banach 空间，则赋弱化拓扑的 E 是拟完备的但不是完备的（II, 第 51 页, 命题 9）。*

拟完备空间是半完备的，因为每个 Cauchy 序列都含于一个闭且有界的子集中（III, 第 3 页, 推论与命题 1）。特别地，可度量化的拟完备局部凸空间是完备的。

在 Hausdorff 拟完备空间中，预紧子集的闭包与闭均衡凸包都是紧的；事实上，它们是预紧的（II, 第 25 页, 命题 3），并且因闭且有界而是完备的（III, 第 3 页, 命题 2）。

命题 9. --- (i) 拟完备局部凸空间的一个闭向量子空间是拟完备的。

\begin{enumerate}
\def\labelenumi{(\roman{enumi})}
\setcounter{enumi}{1}
\item
  拟完备局部凸空间的乘积是拟完备的。
\item
  拟完备局部凸空间的拓扑直和是拟完备的。
\item
  作为闭拟完备子空间的一个序列的严格归纳极限的局部凸空间是拟完备的。
\end{enumerate}

断言 (i) 由注 2（III, 第 2 页）得出，(ii) 由 III, 第 4 页, 推论 2 得出，(iii) 由命题 5（III, 第 5 页）得出，(iv) 由命题 6（III, 第 5 页）得出。

我们以后将看到，拟完备局部凸空间关于一个闭向量空间的商空间未必是拟完备的（IV, 第 63 页, 习题 10）。

命题 10. --- 设 E 是一个局部凸空间，M 是 E 的一个向量子空间，使得 E 的每个点都位于 M 的某个有界子集的闭包中。则从 M 到一个 Hausdorff 拟完备局部凸空间 F 中的每个连续线性映射 f 都唯一扩张为从 E 到 F 中的一个连续线性映射。

该假设蕴含 M 在 E 中稠密，于是 f 唯一扩张为从 E 到 F 的完备化 \(\hat{F}\) 中的一个连续线性映射 f（GT, III, § 3, No.~4, 推论）。但每个 \(x \in F\) 都位于 M 的某个有界子集 B 的闭包中；因此 \(\hat{f}(x)\) 位于 \(f(\mathbf{B})\) 在 \(\hat{F}\) 中的闭包内。而 \(f(\mathbf{B})\) 在 F 中有界，故它在 F 中的闭包是完备的，并且与它在 \(\hat{F}\) 中的闭包相一致。这就证明了 \(\hat{f}(x) \in \mathbf{F}\)。

